\section{Localization near a predecessor span}
\label{l:span:section}

Several signals drawn from the posterior of one exact block satisfy
the same linear equations.  Their affine volume records how close
each new signal is to the span of its predecessors.  We use this
volume in two steps.  First, a relative-entropy calculation divides
the information cost of the block's finite message by the number of
posterior draws.  A draw close to the span of its predecessors incurs
an additional cost.  Second, a short discrete description locates
that draw in a smaller ball, and the decrease in radius pays for the
additional cost.

Throughout this section, \(\mu\) is uniform probability on
\(S^{d-1}\), and all logarithms are natural.  The signal and all
auxiliary variables live in standard Borel spaces.  If the conditional
signal law given \(A\) is \(f_A\mu\), write
\[
 \mathcal D(A)=\mathbb E D(f_A\mu\|\mu).
\]
The variable \(A\) is part of the analysis and may contain continuous
information.  The learner retains only the finite state specified in
Section~\ref{l:model}.

\subsection{The exact equal-label input and an additional moment}

For an integer \(1\le r\le d-2\), let \(\gamma_r\) be the law of an
\(r\)-by-\(d\) standard Gaussian matrix.  If \(X\) has full row rank,
write \(\rho_{0,r}(X,y)\) for the density of \(Xs\), \(s\sim\mu\), and
\(\mu_{X,y}\) for its conditional law.  Put
\[
 X^\dagger=X^{\mathsf T}(XX^{\mathsf T})^{-1}.
\]
On the open ellipsoid \(\|X^\dagger y\|<1\), the conditional law is
uniform on the sphere centered at \(X^\dagger y\), of radius
\(\sqrt{1-\|X^\dagger y\|^2}\), in the affine space parallel to
\(\ker X\).  Its label density is
\begin{equation}
 \rho_{0,r}(X,y)=
 \begin{cases}
 \displaystyle
 \frac{\Gamma(d/2)}{\pi^{r/2}\Gamma((d-r)/2)}
 \frac{(1-\|X^\dagger y\|^2)^{(d-r-2)/2}}
      {\sqrt{\det(XX^{\mathsf T})}},
 &\|X^\dagger y\|<1,\\
 0,&\text{otherwise}.
 \end{cases}
 \label{l:span:base-density}
\end{equation}
One way to verify the exponent is to parameterize the unit sphere in
orthogonal row and null coordinates by
\((u,\sqrt{1-\|u\|^2}\,\theta)\).  Scaling the null-space sphere
contributes \((1-\|u\|^2)^{(d-r-1)/2}\), and the Gram determinant of
the \(u\)-derivatives contributes \((1-\|u\|^2)^{-1/2}\).
The two tangent groups are orthogonal.  Normalizing this density and
changing row coordinates to \(y\) gives
\eqref{l:span:base-density}.  A jointly measurable fiber kernel can
be obtained by projecting an extra Gaussian onto \(\ker X\) and
normalizing it; fixed values are assigned on unused fibers.  On
rank-deficient matrices set the density to zero and use a fixed
fiber probability.

For \(p\ge2\) and \(\mathbf s=(s_1,\ldots,s_p)\), set
\[
 B_{\mathbf s}=[s_2-s_1,\ldots,s_p-s_1],\qquad
 F_{\mathbf s}=\operatorname{col}(B_{\mathbf s}),\qquad
 J(\mathbf s)=\sqrt{\det(B_{\mathbf s}^{\mathsf T}B_{\mathbf s})}.
\]
When \(B_{\mathbf s}\) has full column rank, let
\(\gamma_{r,B_{\mathbf s}}\) denote the probability law of a matrix
whose rows are independent standard Gaussians in \(F_{\mathbf s}^\perp\).

The geometric input is the finite-measure theorem for equal labels
in \cite[Theorem~3.3]{OpenAIReplica2026}.  We use its density-one
spherical instance, stated here in our notation.
\begin{proposition}[Finite measure for equal labels]
\label{l:span:equal-label-assumption}
Let \(r\ge1\), \(p\ge2\), and \(r+p-1\le d-2\).  For every
nonnegative Borel \(H\),
\begin{align}
 &\int\gamma_r(dX)\int_{\mathbb R^r}\rho_{0,r}(X,y)^p
       \int H(X,y,\mathbf s)\,\mu_{X,y}^{\otimes p}(d\mathbf s)\,dy
       \notag\\
 &\quad=(2\pi)^{-r(p-1)/2}
       \int J(\mathbf s)^{-r}\,\mu^{\otimes p}(d\mathbf s)
       \int H(X,Xs_1,\mathbf s)\,\gamma_{r,B_{\mathbf s}}(dX).
 \label{l:span:equal-label-input}
\end{align}
These are finite measures before inserting \(H\).  They are supported
on \(Xs_1=\cdots=Xs_p=y\), full ranks of \(X\) and \(B_{\mathbf s}\),
and positive-radius spherical fibers.  The input uses the exact
density and fiber versions above; its ball-average and Gaussian
approximation versions agree at the constrained label outside a null
set for the finite measure.  The more general theorem permits bounded
nonnegative densities over spherical or Gaussian base probability.
Only the density-one spherical instance is used below.
\end{proposition}
\begin{proof}
Apply \cite[Theorem~3.3]{OpenAIReplica2026} with row count \(r\),
difference count \(p-1\), tuple size \(p\), and base and reweighted
measure both equal to \(\mu\).  Its dimension condition is exactly
\(r+p-1\le d-2\), and its determinant and constrained Gaussian
law are \(J\) and \(\gamma_{r,B_{\mathbf s}}\) above.
\end{proof}

We shall need one more inverse power than appears in the \(r=k\)
identity.  The strict margin in the next lemma makes that use explicit.

\begin{lemma}[The additional inverse-volume moment]
\label{l:span:extra-moment}
Let \(k\ge1\), \(p\ge2\), and \(k+p+1<d\).  Then
\[
 \int J^{-(k+1)}\,d\mu^{\otimes p}<\infty.
\]
\end{lemma}
\begin{proof}
A new sphere point avoids the proper affine span of its predecessors
almost surely, so \(J>0\) for \(\mu^{\otimes p}\)-almost every tuple.
Apply \eqref{l:span:equal-label-input} separately with \(r=k+1\).
Its dimension condition is
\[
 r+(p-1)=k+p\le d-2,
\]
which follows from the integer inequality \(k+p+1<d\).
We check directly that the total mass on the label side is finite.
The exponent in \eqref{l:span:base-density} is
\((d-k-3)/2\ge0\), and hence
\[
 \int \rho_{0,k+1}(X,y)^p\,dy
 \le c_{d,k+1}^{\,p-1}
       \det(XX^{\mathsf T})^{-(p-1)/2}.
\]
Here \(c_{d,k+1}\) is the positive normalizing constant in
\eqref{l:span:base-density}.  Gram--Schmidt on the Gaussian rows
expresses \(\sqrt{\det(XX^{\mathsf T})}\) as the product of their
successive perpendicular lengths.  Conditional on the preceding
rows, these lengths have chi laws in dimensions
\(d,d-1,\ldots,d-k\).  A chi variable in dimension \(v\) has a finite
inverse moment of order \(p-1\) whenever \(v>p-1\), as is seen by
integrating its density \(c_v t^{v-1}e^{-t^2/2}\) at zero and infinity.
The smallest dimension here is \(d-k>p-1\), so iterated conditioning
makes the displayed determinant bound integrable.

Set \(H=1\) in the \(r=k+1\) instance.  Its right side is the positive
constant \((2\pi)^{-(k+1)(p-1)/2}\) times
\(\int J^{-(k+1)}d\mu^{\otimes p}\), because the constrained matrix
law is a probability.  The finite label-side mass proves the claim.
This argument is separate from the later logarithmic estimates.
\end{proof}

\subsection{An unbounded finite-entropy prior}

Suppose \(A\) contains the state at the beginning of a block.  Assume
that the conditional signal law \(f_A\mu\) is supported in an
\(A\)-measurable closed ball \(B(c_A,r_A)\), where \(0<r_A\le1\), and
\begin{equation}
 \mathcal D(A)<\infty,\qquad
 \mathbb E|\log r_A|<\infty.
 \label{l:span:incoming}
\end{equation}
There is no boundedness assumption on \(f_A\).  To obtain a jointly
measurable version, take the Radon--Nikodym derivative of the joint
law of \((A,S)\) with respect to \(\mathcal L(A)\otimes\mu\);
the required absolute continuity follows by integrating the assumed
conditional absolute continuity over \(A\).  Disintegration shows
that this derivative is \(f_A\) for almost every \(A\).  Replacing
the derivative by zero on its infinite-value set gives a finite-valued
Borel version without changing those conditional laws.  All
conditional formulas below are asserted only for this full-measure
set of values of \(A\); arbitrary values may be used elsewhere.

Fix \(1\le k\) with \(k+p+1<d\).  Draw \(X\sim\gamma_k\)
independently of \((A,S)\), put \(Y=XS\), and draw the block destination
\(w\) by the learner's kernel of the incoming state and \((X,Y)\).
Assume that \(w\) has at most \(W\) values.  Given \((A,X,Y,w)\), draw
\(s_1,\ldots,s_p\) independently from the conditional law of \(S\)
given \((A,X,Y)\).  The destination provides no further information
once these data are known.  Thus every \((s_i,A,X,Y,w)\) has the
original \((S,A,X,Y,w)\) law.

We will replace \(S\) by a signal having the same joint law with
\((A,w)\), and enlarge \((A,w)\) to an analysis variable \(A'\)
that determines a supporting ball \(B(c',r')\), with
\(0<r'\le r_A\).  For any fixed coefficient \(a\ge k+p\), our
one-block objective is
\[
 \mathcal D(A')+a\mathbb E\log r'
 \le \mathcal D(A)+a\mathbb E\log r_A
       +\frac{\log W}{p}+\frac{k}{2}+Cp+2a\log2,
\]
where \(C\) is absolute.  We obtain this bound by retaining a
uniformly selected posterior draw and its predecessors, then
discretely locating that draw near their affine span.  First we
derive the exact likelihood of the tuple and prove its logarithms
integrable, so that the entropy chain rule can be applied.

For fixed \(A\) and full-row-rank \(X\), define
\[
 \rho_A(X,y)=
 \begin{cases}
 \displaystyle \rho_{0,k}(X,y)\int f_A(s)\,\mu_{X,y}(ds),
 &\rho_{0,k}(X,y)>0\ \text{and }\displaystyle\int f_A\,d\mu_{X,y}<\infty,\\
 0,&\text{otherwise}.
 \end{cases}
\]
Set \(\rho_A=0\) on rank failures.  Kernel integration gives a
jointly measurable function of
\((A,X,y)\).  Disintegration shows that it is a conditional label
density; the zero assignment on infinite fiber means changes it
only on a Lebesgue-null set for almost every \((A,X)\).  The
piecewise definition also avoids a product \(0\cdot\infty\) on
unused fibers.  The actual label has \(0<\rho_A(X,Y)<\infty\)
almost surely, and on that set its posterior is
\[
 \frac{\rho_{0,k}(X,Y)}{\rho_A(X,Y)}f_A\,\mu_{X,Y}.
\]
This definition fixes the version used below.

Let \(\pi_A\) be the law of \((\mathbf s,X)\) in this experiment,
with \(w\) discarded.  Define the probability measure
\[
 \lambda(d\mathbf s,dX)
   =\mu^{\otimes p}(d\mathbf s)\,
      \gamma_{k,B_{\mathbf s}}(dX),
\]
with an arbitrary definition on the \(\mu^{\otimes p}\)-null set
where \(J=0\).

\begin{lemma}[Replica likelihood with finite logarithms]
\label{l:span:likelihood}
Under \eqref{l:span:incoming},
\begin{equation}
 \frac{d\pi_A}{d\lambda}(\mathbf s,X)
  =(2\pi)^{-k(p-1)/2}J(\mathbf s)^{-k}
       \frac{\prod_{i=1}^p f_A(s_i)}
            {\rho_A(X,Xs_1)^{p-1}}.
 \label{l:span:likelihood-formula}
\end{equation}
The multiplier is defined as zero unless \(J>0\) and its
denominator is finite and positive.  Under the replica experiment,
each \(\log f_A(s_i)\), \(\log\rho_A(X,Y)\), and \(\log J\) is
absolutely integrable after averaging over \(A\).
\end{lemma}
\begin{proof}
For each such value of \(A\), on the label side of
\eqref{l:span:equal-label-input} with \(r=k\), insert the
everywhere-defined nonnegative Borel multiplier
\[
 \begin{cases}
 \displaystyle\frac{\prod_i f_A(s_i)}{\rho_A(X,y)^{p-1}},
       &0<\rho_A(X,y)<\infty,\\
 0,&\text{otherwise}.
 \end{cases}
\]
The identity applies first to bounded truncations and then to this
multiplier by monotone convergence.  The label-side measure becomes
the actual label law \(\rho_A(X,y)\,dy\) followed by \(p\) normalized
posterior draws.  The other side is the right side of
\eqref{l:span:likelihood-formula} relative to \(\lambda\).
This proves the likelihood formula.  In particular, no bounded
posterior theorem has been applied to \(f_A\); the only geometric
input was the identity for \(\mu\).  Values on invalid label fibers
are handled by the multiplier itself, so no unproved transfer of a
Lebesgue-null set to the singular reference is being used.

We justify logarithmic integrability before taking any entropy
difference.  Since \(s_i\mid A\) has density \(f_A\),
\[
 \int f_A(\log f_A)_-\,d\mu\le 1/e.
\]
Together with \(\mathcal D(A)<\infty\), this gives averaged absolute
integrability of \(\log f_A(s_i)\).  Let \(g\) be standard Gaussian
density on \(\mathbb R^k\).  For almost every \((A,X)\),
\[
 \int\rho_A(X,y)
       \left(\log\frac{g(y)}{\rho_A(X,y)}\right)_+dy
 \le \frac1e\int g(y)\,dy=\frac1e.
\]
Moreover \(\mathbb E\|Y\|^2=k\), because \(X\) is fresh and
\(\|S\|=1\).  Since
\(-\log g(Y)=\tfrac{k}{2}\log(2\pi)+\tfrac12\|Y\|^2\),
the negative part of \(\log\rho_A(X,Y)\) is integrable.

Put \(c=(2\pi)^{-k(p-1)/2}\) and
\[
 Z=c\int J^{-k}\,d\mu^{\otimes p}.
\]
The \(r=k\) input shows \(0<Z<\infty\).  Let \(Q\) be the probability
with density \(cJ^{-k}/Z\) relative to \(\lambda\).  The likelihood
formula gives \(\pi_A\ll\lambda\).  Since \(0<J<\infty\)
\(\lambda\)-almost surely, the density of \(Q\) is finite and strictly
positive there; thus \(Q\) and \(\lambda\) have the same null sets,
and \(\pi_A\ll Q\).  We may therefore write
\begin{equation}
 \log\frac{d\pi_A}{dQ}
   =\log Z+\sum_{i=1}^p\log f_A(s_i)
       -(p-1)\log\rho_A(X,Y).
 \label{l:span:tilted-log}
\end{equation}
Its positive part satisfies
\[
 \left(\log\frac{d\pi_A}{dQ}\right)_+
 \le |\log Z|+\sum_i(\log f_A(s_i))_+
       +(p-1)(\log\rho_A(X,Y))_-,
\]
so it is integrable by the two bounds just proved, without using
\(\log J\).  Its negative part has integral at most \(1/e\):
if \(v=d\pi_A/dQ\), then
\(\int v(\log v)_-\,dQ\le1/e\).  Hence
\(\mathbb E D(\pi_A\|Q)<\infty\), and the log likelihood in
\eqref{l:span:tilted-log} is absolutely integrable on average.
Rearranging that identity now proves full integrability of
\(\log\rho_A(X,Y)\).  This step has not used \(\log J\).

By Lemma~\ref{l:span:extra-moment},
\[
 \int e^{(\log J)_-}\,dQ
 \le \frac cZ\left(
       \int J^{-k}\,d\mu^{\otimes p}
       +\int J^{-(k+1)}\,d\mu^{\otimes p}\right)<\infty.
\]
For probabilities \(P,Q\) and nonnegative \(F\), exponential tilting
gives
\[
 \mathbb E_P F\le D(P\|Q)+\log\mathbb E_Q e^F
\]
first for bounded \(F\), and then for nonnegative \(F\) by truncation.
This is the bounded variational inequality of
\cite[Equations~(1.1)--(1.2)]{DupuisMao2022}, with the limiting step
given here.  Apply it to \(F=(\log J)_-\) conditionally on \(A\) and
average.  Its right side is finite.  The positive part of \(\log J\)
is at most \((p-1)\log2\) by Hadamard's inequality, since every chord
has length at most two.  This completes the integrability proof.
\end{proof}

\subsection{Information in one selected draw}

For \(i\ge2\), let
\[
 \delta_i=\operatorname{dist}\!
   \left(s_i,\operatorname{aff}(s_1,\ldots,s_{i-1})\right),
 \qquad
 \Delta_i=\max\{0,\log(r_A/\delta_i)\},
 \qquad \Delta_1=0.
\]
Gram--Schmidt gives \(J=\prod_{i=2}^p\delta_i\).  The likelihood
formula implies absolute continuity of the tuple marginal with
respect to \(\mu^{\otimes p}\), so every \(\delta_i>0\) almost surely.
Also \(\delta_i\le2\) and
\(\Delta_i\le(-\log\delta_i)_+\).  Thus
\[
 (-\log\delta_i)_+
 \le(-\log J)_++(p-2)\log2,
\]
and Lemma~\ref{l:span:likelihood} implies
\(\mathbb E\Delta_i<\infty\).

Choose an independent uniform \(I\in\{1,\ldots,p\}\), use \(s_I\)
as the signal for the next block, and set
\[
 A_*=(A,w,I,s_1,\ldots,s_{I-1}).
\]
The triple \((s_I,A,w)\) has the original \((S,A,w)\) law, since
each draw has that marginal and \(I\) is independent.  Uniform
selection turns the entropy-chain average into information about
the selected draw; the following bound is not asserted for every
fixed draw.

\begin{lemma}[Information of a selected draw]
\label{l:span:selected}
The conditional signal law given \(A_*\) is absolutely continuous
with respect to \(\mu\), and
\begin{equation}
 \mathcal D(A_*)\le \mathcal D(A)+\frac{\log W}{p}
                    +\frac{k}{2}+k\,\mathbb E\Delta_I.
 \label{l:span:selected-bound}
\end{equation}
\end{lemma}
\begin{proof}
Taking logarithms in \eqref{l:span:likelihood-formula} is legitimate
by Lemma~\ref{l:span:likelihood}.  Discard \(X\) by data processing,
then condition on \(w\).  The latter step adds the conditional
mutual information of the tuple and \(w\), at most \(\log W\).
The product-reference chain rule, in the finite form discussed in
\cite[Section~3.3, Equations~(13)--(18)]{Austin2020}, therefore yields
\begin{align}
 \mathbb E D(\mathcal L(\mathbf s\mid A,w)\|\mu^{\otimes p})
 &\le \log W+p\mathcal D(A)
       +(p-1)\mathbb E[-\log\rho_A(X,Y)] \notag\\
 &\quad-\frac{k(p-1)}2\log(2\pi)-k\mathbb E\log J.
 \label{l:span:tuple-kl}
\end{align}
The label entropy is bounded by its cross entropy with
\(N(Xc_A,r_A^2I_k)\).  Freshness and the containing ball give
\(\mathbb E(\|Y-Xc_A\|^2/r_A^2\mid A)\le k\), whence
\[
 \mathbb E[-\log\rho_A(X,Y)]
 \le \frac{k}{2}\log(2\pi)+k\mathbb E\log r_A+\frac{k}{2}.
\]
All terms are finite by \eqref{l:span:incoming} and the preceding
lemma.  Substitute this estimate and
\(\log J=\sum_{i=2}^p\log\delta_i\) into
\eqref{l:span:tuple-kl}.  Since
\(\log(r_A/\delta_i)\le\Delta_i\), the result is
\[
 \mathbb E D(\mathcal L(\mathbf s\mid A,w)\|\mu^{\otimes p})
 \le \log W+p\mathcal D(A)+\frac{k(p-1)}2
            +k\mathbb E\sum_{i=1}^p\Delta_i.
\]
This finite relative entropy also implies that, almost surely,
each conditional coordinate law given its predecessors is
absolutely continuous with respect to \(\mu\).  The product-reference
chain rule and the uniform choice of \(I\) give the exact identity
\[
 \mathcal D(A_*)
 =\frac1p\,\mathbb E D\!\left(
     \mathcal L(\mathbf s\mid A,w)\middle\|\mu^{\otimes p}\right).
\]
Dividing the preceding bound by \(p\) and using
\((p-1)/p\le1\) gives
\eqref{l:span:selected-bound}.
\end{proof}

\subsection{A radius credit from the predecessor span}

We next encode a smaller ball when the selected draw is close to its
predecessor span.  The code is discrete even though the predecessor
coordinates, already in \(A_*\), are continuous.

\begin{lemma}[A discrete description of a smaller ball]
\label{l:span:localization}
There is a discrete code \(C\), determined by \((A_*,s_I)\), and an
\(A'=(A_*,C)\)-measurable ball \(B(c',r')\) supporting the conditional
signal law, for which
\begin{align}
 0<r'&\le r_A,&
 \log r'&\le\log r_A-\Delta_I+2\log2,
 \label{l:span:radius}\\
 H(C\mid A_*)&\le p\,\mathbb E\Delta_I+C_1p,&
 \mathcal D(A')&\le\mathcal D(A_*)+H(C\mid A_*).
 \label{l:span:code}
\end{align}
Here \(C_1\) is absolute.  Both \(\mathcal D(A')\) and
\(\mathbb E|\log r'|\) are finite.
\end{lemma}
\begin{proof}
If \(I=1\), set \(j=0\).  For \(I\ge2\), set
\[
 j=\max\{0,\lfloor\log_2(r_A/\delta_I)\rfloor\}.
\]
When \(j=0\), keep the old ball.  When \(j\ge1\), the predecessor
span has dimension \(v=I-2\) almost surely.  Use orthonormal
coordinates in that affine span based at \(s_1\).  The projection
of \(s_I-s_1\) has norm at most \(2r_A\), because both points lie
in the old ball.

For every integer pair \((j,v)\), fix an ordered \(2^{-j}\)-net of
the radius-two ball in \(\mathbb R^v\) with at most
\((1+4\cdot2^j)^v\) points.  A maximal \(2^{-j}\)-separated set has
this size by comparing the volumes of its disjoint radius-\(2^{-j-1}\)
balls with the containing ball.  In dimension zero use one point.
Record the first net point within \(2^{-j}\) of the projection
coordinates divided by \(r_A\).  After scaling and translating back,
let \(c'\) be that point of the predecessor span.  The in-span error
and perpendicular error are both at most \(r_A2^{-j}\).  Hence
\(B(c',2r_A2^{-j})\) contains \(s_I\); set \(r'=2r_A2^{-j}\).
It is at most \(r_A\) for \(j\ge1\).

Figure~\ref{l:span:net-figure} depicts this step, writing
\(\eta=r_A2^{-j}\) and \(F=\operatorname{aff}(s_1,\ldots,s_{I-1})\)
within the illustration. Only the \(v\)-dimensional projection must be
encoded; the perpendicular error is absorbed by the decoded radius.
The radius does not increase, but the balls need not be nested.

\begin{figure}[htbp]
\centering
\begin{tikzpicture}[x=1.12cm,y=1.12cm,>=Latex,
  every node/.style={font=\small}]
  \fill[blue!4] (1.6,-.4) circle[radius=1.3];
  \draw[gray,dashed] (-.5,.2) circle[radius=2.6];
  \draw[blue!70!black,thick] (1.6,-.4) circle[radius=1.3];
  \draw[gray] (-3.3,-.4) -- (3.1,-.4);
  \node[anchor=west] at (3.12,-.4) {$F$};
  \node[anchor=west,align=left] at (3.7,.9)
    {$F=\operatorname{aff}(s_1,\ldots,s_{I-1})$\\
     $\dim F=v=I-2$};
  \fill (-2.1,-.4) circle[radius=.035];
  \node[below=4pt] at (-2.1,-.4) {$s_1$};
  \fill (-.7,-.4) circle[radius=.035];
  \node[below=4pt] at (-.7,-.4) {$s_2$};
  \fill (1.1,.12) circle[radius=.045];
  \node[above left=2pt] at (1.1,.12) {$s_I$};
  \draw[densely dashed] (1.1,.12) -- (1.1,-.4);
  \draw (1.1,-.29) -- (1.21,-.29) -- (1.21,-.4);
  \fill (1.1,-.4) circle[radius=.03];
  \node[below left=3pt] at (1.1,-.4) {$P$};
  \draw[<->] (.85,.12) -- (.85,-.4);
  \node[anchor=east,fill=white,inner sep=1.5pt]
    at (.77,-.12) {$\delta_I\le\eta$};
  \fill[blue!70!black] (1.6,-.4) circle[radius=.05];
  \node[below right=3pt,blue!70!black] at (1.6,-.4) {$c'$};
  \draw[<->] (1.1,-.94) -- (1.6,-.94);
  \node[below=2pt] at (1.35,-.94) {$\le\eta$};
  \draw[blue!70!black,->] (1.6,-.4) -- (2.54,.5);
  \node[anchor=west,blue!70!black,fill=white,inner sep=1.5pt]
    at (2.25,.05) {$r'=2\eta$};
  \draw[gray,->] (-.5,.2) -- (-1.99,2.32);
  \fill[gray] (-.5,.2) circle[radius=.03];
  \node[above left=2pt,gray,fill=white,inner sep=1.5pt]
    at (-.5,.2) {$c_A$};
  \node[anchor=east,gray] at (-1.36,1.36) {$r_A$};
  \node[anchor=west,align=left] at (3.7,.2)
    {$\eta=r_A2^{-j}$, $j\ge1$};
  \begin{scope}[shift={(3.8,-.9)},x=.58cm,y=.58cm]
    \draw[gray] (0,0) -- (5.4,0);
    \draw[gray,densely dotted] (-.5,0) -- (0,0);
    \draw[gray,densely dotted] (5.4,0) -- (5.9,0);
    \draw[blue!55,thick] (0,.10) -- (1.75,.10);
    \draw[blue!55,thick] (.92,.18) -- (2.92,.18);
    \draw[blue!55,thick] (2.52,.10) -- (4.52,.10);
    \draw[blue!55,thick] (3.97,.18) -- (5.4,.18);
    \foreach \x in {.75,1.92,3.52,4.97}
      \fill[blue!70!black] (\x,0) circle[radius=.065];
    \fill (1.42,0) circle[radius=.05];
    \node[below=3pt] at (1.42,0) {$P$};
    \node[below right=3pt,blue!70!black] at (1.92,0) {$c'$};
    \draw[<->] (1.92,.64) -- (2.92,.64);
    \node[above=1pt] at (2.42,.64) {$\eta$};
  \end{scope}
  \node[anchor=north west,align=left] at (3.7,-1.45)
    {One-dimensional covering net\\[4pt]
     $\#\text{net}\le(1+4\cdot2^j)^v$\\[4pt]
     Scale and index: $(v+1)j+O(v+1)$ bits};
\end{tikzpicture}
\caption{Predecessor-span localization, shown schematically in two
dimensions. Only the projection of \(s_I\) onto \(F\) is quantized:
the projected point \(P\) lies in the radius-\(2r_A\) ball about
\(s_1\) in the \(v\)-dimensional span. A fixed \(\eta\)-cover supplies
the first qualifying ordered net point \(c'\); the inset illustrates
part of such a cover in one dimension, with omitted stretches dotted.
Both errors are at most \(\eta\), so \(B(c',2\eta)\) contains \(s_I\).
The scale and net index have the displayed code-length bound, with
dimension \(v\). The new radius does not increase, but the new ball
need not lie inside the old one.}
\label{l:span:net-figure}
\end{figure}

For \(j\ge1\),
\(j\log2\le\Delta_I<(j+1)\log2\).  Thus
\[
 \log r_A-\Delta_I\le \log r'
       \le\log r_A-\Delta_I+2\log2.
\]
For \(j=0\), either \(\Delta_I=0\) or \(\Delta_I<\log2\), and the
same upper bound holds with \(r'=r_A\).  These inequalities give
\eqref{l:span:radius} and, using the finite expectations already
proved, \(\mathbb E|\log r'|<\infty\).

Encode \(j\) by the unary prefix code of length \(j+1\) bits.
For \(j\ge1\), append a fixed-length net index of
\(\lceil v\log_2(1+4\cdot2^j)\rceil\) bits.  Given \(A_*\), the
dimension \(v\) is known, so the combined code is prefix-free.
Since \(1+4\cdot2^j\le5\cdot2^j\), its length is at most
\((1+v)j+C v+2\).  Now \(1+v\le p\) and
\(j\log2\le\Delta_I\), so its expected length in nats is at most
\(p\mathbb E\Delta_I+C_1p\).  We give the countable conditional
entropy step explicitly.  For a conditioning value \(a_*\) with finite
conditional expected length, let \(\ell_{a_*}(c)\) be the bit length
and let \(p_{a_*}(c)\) be the conditional code distribution.  Kraft's
inequality \cite{Kraft1949}, in the finite or countable form stated in
\cite[Section~2.4.2, Theorem~2.4.2]{Duchi2023}, gives
\[
 0<K_{a_*}:=\sum_c2^{-\ell_{a_*}(c)}\le1.
\]
Normalize \(q_{a_*}(c)=2^{-\ell_{a_*}(c)}/K_{a_*}\).  Its cross entropy is
\[
 \sum_c p_{a_*}(c)\log\frac1{q_{a_*}(c)}
 =(\log2)\mathbb E[\ell_{a_*}(C)\mid a_*]+\log K_{a_*}
 \le(\log2)\mathbb E[\ell_{a_*}(C)\mid a_*]<\infty.
\]
For a finite partition of the code alphabet with one remaining tail
atom, nonnegative relative entropy bounds the partition entropy by
its \(q_{a_*}\) cross entropy, and that cross entropy is at most the
displayed full cross entropy.  Refining such partitions gives the
countable entropy by monotone convergence.  Hence
\(H(C\mid a_*)\le(\log2)\mathbb E[\ell_{a_*}(C)\mid a_*]\).
Finite conditional expected length holds for almost every \(a_*\)
because its average is finite.  Averaging proves the first inequality
in \eqref{l:span:code}.

Gram--Schmidt with fixed conventions supplies measurable
orthonormal coordinates from the predecessor differences, and the
fixed ordering resolves net ties.  Define the code arbitrarily on
the null degeneracy set.  The decoded ball is therefore measurable.
The pointwise construction puts \(s_I\) in that decoded ball outside
the null degeneracy set.  The event
\(\{\|s_I-c'(A')\|\le r'(A')\}\) is Borel and has probability one.
Disintegrating it over \(A'\) shows that its conditional probability
is one for almost every \(A'\).  Thus the decoded ball supports the
conditional signal law.
Conditioning an absolutely continuous law on a discrete code keeps
absolute continuity almost surely.  The conditional entropy chain
rule gives the second inequality in \eqref{l:span:code}, and its
finite right side proves the remaining finiteness claim.
\end{proof}

\subsection{The sample bound}

The preceding two lemmas now give a complete potential argument.
The result uses the local model and the local cap and full-data
calculation of Section~\ref{l:model}; the only companion input remains
\eqref{l:span:equal-label-input}.

\begin{proposition}[A predecessor-span block bound]
\label{l:span:block-bound}
For sufficiently large \(d\), put \(K=p=\lfloor d/8\rfloor\).
Consider a learner of Section~\ref{l:model} with horizon \(T\) and
\(M\) persistent bits, and set \(W=(T+2)2^M\).  If
\(\log W\le2d^2\), then its uniform-sphere angular success
probability \(\pi\) at \(0<\epsilon\le1/10\) satisfies
\begin{equation}
 \pi(d-1)\log(1/\epsilon)
 \le C d\left(\left\lceil T/K\right\rceil+1\right)
 \label{l:span:success-bound}
\end{equation}
for an absolute \(C\).  In particular, success at least \(2/3\)
implies
\[
 \left\lceil T/K\right\rceil
       \ge c\log(1/\epsilon)-C'
\]
with absolute \(c>0,C'>0\).
\end{proposition}
\begin{proof}
Condition on the shared rule randomness.  For almost every value, the
joint initialization independence in Section~\ref{l:model} leaves
the initial state independent of the signal and all sample rows.
By Lemmas~\ref{l:borel} and \ref{l:prior-ac}, the conditional
uniform-prior experiment has Borel rule versions.  We prove the
displayed numerical bound for each such value, with its conditional
success probability.  Averaging that bound at the end handles shared
randomness without requiring jointly selected Borel versions.

Pad an early-stopping program to \(T\) observations by retaining its
terminal state and stopping index, and then ignoring the remaining
pairs.  At each index the padded state has at most \(W\) values, and
the final output law uses only that state.  Divide the observations
into blocks of length \(K\), allowing a shorter final block.  Set
\[
 a=K+p,\qquad
 \Phi(A,r_A)=\mathcal D(A)+a\mathbb E\log r_A.
\]
For every nonempty block length \(1\le k\le K\), the inequalities
\(k+p+1<d\) and \(a\le(d-1)/2\) hold once \(d\) is sufficiently
large.  Initially \(A\) is the signal-independent initial state and
the containing ball is \(B(0,1)\), so \(\Phi=0\).

At a block, select a replica and encode its new ball as above.
Combining \eqref{l:span:selected-bound}--\eqref{l:span:code} gives
\begin{align}
 \Phi(A',r')
 &\le\Phi(A,r_A)+\frac{\log W}{p}+\frac{k}{2}+C_1p+2a\log2
           +(k+p-a)\mathbb E\Delta_I \notag\\
 &\le\Phi(A,r_A)+C_2d.
 \label{l:span:drift}
\end{align}
Indeed \(k+p-a\le0\), while \(p\ge d/16\) and
\(\log W\le2d^2\) make the width charge \(O(d)\).
The cancellation has a direct explanation: the likelihood
charges \(k\Delta_I\), the code charges at most \(p\Delta_I\), and
the radius credit is \(a\Delta_I=(K+p)\Delta_I\).

Every posterior draw has the correct joint marginal with the block
destination.  Uniformly selecting one draw therefore preserves the
learner's state--signal marginal, and recording its code changes no
marginal law.  Generate the next Gaussian block independently of the
pair consisting of the current signal and augmented analysis
variables.  Since the learner's transition uses
the past only through its state, induction preserves the actual
state--signal marginal at every boundary.  The predecessor points
and codes are never given to the learner.  All required finiteness
properties propagate by Lemma~\ref{l:span:localization}.  Thus after
\(b=\lceil T/K\rceil\) blocks,
\[
 \Phi\le C_2db.
\]
For \(T=0\) this holds directly with \(b=0\).

Write \(n=d-1\).  Lemma~\ref{l:caps-endpoint} gives
\(\mu(B(c_A,r_A))\le(2r_A)^n\).  A density supported in a set \(E\)
has relative entropy at least \(-\log\mu(E)\), by comparison with
\(\mu(\,\cdot\mid E)\).  Therefore
\[
 D(f_A\mu\|\mu)\ge-n\log(2r_A)
\]
almost surely, and
\begin{equation}
 \Phi\ge\left(1-\frac an\right)\mathcal D(A)-a\log2
       \ge\frac12\mathcal D(A)-a\log2.
 \label{l:span:coercivity}
\end{equation}
The unconditional signal law remains \(\mu\), so
\(\mathcal D(A)=I(S;A)\).  The final \(A\) includes the padded state
\(w\), and hence \(\mathcal D(A)\ge I(S;w)\).

Under the product of \(\mu\) and the output marginal, angular
success has probability at most \(\epsilon^n\), by
Lemma~\ref{l:caps-endpoint}.  Data processing through the terminal
output and then through the success event gives
\[
 I(S;w)\ge I(S;\widehat S)
       \ge \pi n\log(1/\epsilon)-\log2
\]
by \eqref{l:binary}.  This holds with the success probability for
the fixed shared seed.  Combining it with the potential upper bound
and \eqref{l:span:coercivity} proves
\eqref{l:span:success-bound} with constants independent of that
seed.  Averaging proves the asserted bound for the original shared
randomness as well.
\end{proof}

\begin{theorem}[The full accuracy range by predecessor-span localization]
\label{l:span:theorem}
Let \(M(d)=o(d^2)\) be integer-valued and let
\(0<\epsilon(d)\le1/10\).  A family of learners in
Section~\ref{l:model} whose angular success probability is at least
\(2/3\) for a uniform spherical signal must use
\[
 T(d)\ge c\,d\log(1/\epsilon(d))
\]
for all sufficiently large \(d\), where \(c>0\) is absolute.
The eventual threshold may depend on the memory sequence.  The
same conclusion follows from success at least \(2/3\) for every
fixed signal.
\end{theorem}
\begin{proof}
Put \(L=\log(1/\epsilon)\ge\log10\).  The last statement follows by
averaging over the uniform signal, so it suffices to prove the average
statement.  There are at most \((T+1)2^M\) terminal time--state
pairs.  For each fixed pair its fresh output law succeeds on uniform
mass at most \(\epsilon^{d-1}\), by the cap bound and integration
over that output law.  More explicitly, for a terminal pair \(h\)
let \(a_h(s)\) be that fixed output law's success probability at
signal \(s\), and let \(t_h(s)=\Pr(h\text{ is reached}\mid S=s)\).
The pair's contribution is
\[
 \int t_h(s)a_h(s)\,\mu(ds)
 \le\int a_h(s)\,\mu(ds)\le\epsilon^{d-1},
\]
because \(0\le t_h(s)\le1\).  Summing over all pairs, for almost
every shared seed and hence after averaging, yields
\[
 \frac23\le (T+1)2^M\epsilon^{d-1}.
\]
Equivalently,
\begin{equation}
 (d-1)L\le M\log2+\log(T+1)+\log(3/2).
 \label{l:span:output-count}
\end{equation}

At indices where \(T\ge dL\) the conclusion is immediate.  At the
other indices, \eqref{l:span:output-count} gives
\[
 (d-1)L\le M\log2+\log(dL+1)+\log(3/2).
\]
For all sufficiently large \(d\), uniformly for \(L\ge\log10\),
\(\log(dL+1)\le dL/6\), while \((d-1)L\ge dL/2\).
Consequently \(dL\le3M\log2+3\log(3/2)=o(d^2)\) along such
short indices.  In particular \(\log((T+2)2^M)\le2d^2\) eventually,
so Proposition~\ref{l:span:block-bound} applies.  It gives
\[
 T\ge K\{c_0L-C_0-1\}
\]
with absolute \(c_0>0,C_0\).  Above a sufficiently large absolute
threshold for \(L\), use \(K\ge d/16\) and absorb the constants to
obtain \(T\ge c_1dL\).

For completeness, the local full-data argument supplies the remaining
constant-accuracy range with explicit constants.  If \(T\le d/8\),
reveal every row and label.  The conditional unknown component is
uniform on a residual sphere of radius \(R\).  By isotropy and
Markov's inequality,
\[
 \Pr\{R<1/2\}\le \frac{4T}{3d}\le\frac16.
\]
On \(R\ge1/2\), antipodal residual points correspond to signals
separated by at least one, so a radius-\(\epsilon\) error ball with
\(\epsilon\le1/10\) contains at most one point of each pair.
Conditional success is at most \(1/2\), also for a randomized
estimate.  Total success is at most
\[
 \frac16+\frac56\cdot\frac12=\frac7{12}<\frac23.
\]
Thus \(T>d/8\).  This proves \(T\ge c_2dL\) throughout the bounded
range of \(L\) left above.  Taking the minimum of the absolute
constants completes the proof, including nonmonotone accuracy
sequences and early stopping.
\end{proof}
